\section{The algorithm and its finite implementation}
\label{sec:implementation}

We now combine the two modular procedures.  Their roles are complementary:
the checksum counts every subset in a modular bin, whereas the extraction
procedure identifies exactly the terms that do not belong to the integer
bin.  We first describe the resulting decision rule, then supply the finite
arithmetic and bounded sampling needed to prove Theorem~\ref{thm:main} in
the stated machine model.

\subsection{Subtracting aliases}

Choose a fixed sufficiently large integer $n_*$, whose requirements will be
specified below.  If $n<n_*$ or the bit-length guard
\eqref{eq:bit-guard} fails, use exact meet-in-the-middle on the original
instance.  Otherwise choose the tie weights and form $c_i,t_s$ as in
\eqref{eq:transformed}.  Run the preliminary procedure of
Proposition~\ref{prop:preliminary}.  Return its answer if it decides the
instance; its outcome \textsc{Continue} means only that we proceed to the
following steps.

Run the alias procedure of Theorem~\ref{thm:aliases} to choose the large
prime $P$ and obtain its alias lists.  Write $\mathcal A_s$ for its
deduplicated set of aliases for target $t_s$.  Every retained subset is checked using the
original transformed integers, so even on an unsuccessful extraction run
the set $\mathcal A_s$ contains only genuine aliases.  On a successful run
it contains all of them.  Next draw and retain the phases
$\phi_1,\ldots,\phi_n$ used in Theorem~\ref{thm:checksum}, and compute its
estimates $\widehat V_{t_s\bmod P}$ for all target residues.  Repeated
residues can share an estimate.  For each $s$, form the finite-precision
quantity
\begin{equation}
 \widehat B_s
 =\widehat V_{t_s\bmod P}
  -\sum_{I\in\mathcal A_s}
       e_P\left(\sum_{i\in I}\phi_i\right).
 \label{eq:subtract-aliases}
\end{equation}
Each phase in this subtraction is evaluated from the retained integer
residues $\phi_i$; it is not resampled.  Deduplication is by the binary
subset itself, separately for each $s$, not by its sum or phase.

Output \textsc{Yes} if $|\widehat B_s|\ge 1/2$ for some $s$, and
\textsc{No} otherwise.  Since the computed real and imaginary parts are
dyadic rationals, this test is implemented exactly by comparing the sum of
their squares with $1/4$.  Any failure declared by a capped main procedure
or a main-procedure random sampler results in \textsc{No}.  This convention
does not apply to the preliminary outcome \textsc{Continue}.

The cancellation behind this decision rule is exact before numerical
rounding.  When extraction is complete,
\begin{equation}
 V_{t_s\bmod P}
 -\sum_{I\in\mathcal A_s}e_P\left(\sum_{i\in I}\phi_i\right)
 = B_s,
 \qquad
 B_s:=\sum_{\substack{I\subseteq[n]\\\sum_{i\in I}c_i=t_s}}
             e_P\left(\sum_{i\in I}\phi_i\right).
 \label{eq:exact-bin-checksum}
\end{equation}
In particular, an empty exact bin has $B_s=0$, and a bin containing one
subset has $|B_s|=1$, for every choice of phases.  Thus isolation, rather
than an anti-concentration estimate for a many-term sum, is what turns
checksum estimation into a decision algorithm.

\subsection{Finite arithmetic and bounded sampling}
\label{sec:finite-arithmetic}

The preceding sections describe finite combinatorial constructions, but
their analyses use exact uniform draws and exact complex values.  We now
implement both without unbounded running times or real-arithmetic
instructions.  All constants below are fixed independently of the input
and of the exponent $c$ in the polynomial-bit convention.

\paragraph{Integer data and implicit indices.}
The input integers, transformed sums, targets, and signed record weights
have $O(b+\log(n+2))$ bits.  Masks have $n$ bits, and the residues, digit
indices, block-table indices, and their tuples have polynomially many
bits in $n$.  Arithmetic on these objects is performed with ordinary
multiword integer algorithms.  For example, shift-and-add multiplication
and binary long division use a polynomial number of bit operations in
their operand lengths, and each bit operation is simulated by a constant
number of allowed word operations.  Signs are stored separately.

There is one useful distinction between a stored list and its implicit
Cartesian products.  In a sampled checksum pattern let
$E_\pi^{\rm hi},E_\pi^{\rm lo}$ be the numbers of emitted high and low
entries; retain the notation $L_\pi$ for the number of matches.  The entry
caps imply on every execution, not only on the good event, that
\begin{equation}
 L_\pi\le E_\pi^{\rm hi}E_\pi^{\rm lo}
 \le 2^{2\lceil(\tau+20\eps)n\rceil}.
 \label{eq:implicit-match-bitlength}
\end{equation}
Thus the possibly much larger match set is counted and indexed with
$O(n)$-bit integers.  The group counts, prefix sums, and blockwise
multiplicities in \eqref{eq:checksum-match-count} and
\eqref{eq:checksum-multiplicity} therefore require only polynomial
arithmetic per operation.  Those constructions index matches without
an array indexed by all matches or by all $P$ residues.

The list generators in Sections~\ref{sec:checksum} and~\ref{sec:aliases}
emit records one at a time and check their caps before the next write.
The alias exclusion uses binary searches rather than a scan of the
forbidden interval, as proved in Section~\ref{sec:aliases}.
Merge sorting uses at most a constant multiple of the stored records as
temporary storage, with polynomial arithmetic per comparison.
The distinct-sum preliminary procedure
can merge a list with its translate, discarding duplicates while emitting
the next list; its temporary length is at most a constant multiple of the
cap.

\paragraph{Uniform draws with a deterministic time limit.}
A fair bit is obtained, for example, from the low bit of a fresh random
word.  To draw uniformly from $\{0,\ldots,q-1\}$, draw
$\lceil\log q\rceil$ independent bits and reject if their integer value
is at least $q$; for $q=1$ return zero without drawing.  Each proposal is
accepted with probability greater than $1/2$.  Stop with failure after
$n^3$ rejected proposals.  Before that failure this is exactly the usual
uniform sampler, and the probability of failure at any given call is at
most $2^{-n^3}$.  This implements random indices, residues, and a rational
Bernoulli probability $p/q$ by testing a uniform index against $p$.
Uniform subsets of prescribed size are obtained by the first positions
of a permutation generated by successive uniform choices among the
remaining positions.  The independent flip bits need no rejection.

For a prime in $[2^h,2^{h+1})$, draw uniform integers in this interval and
test each for primality, stopping at the first prime.  Try at most $n^3$
candidates.  Conditional on a prime being found, it is uniform among the
primes of the interval.  The prime interval estimate used in
Lemma~\ref{lem:prime-collision} gives an acceptance probability at least
$c_0/n$ for a fixed $c_0>0$ at the two prime scales used here.  Thus failure
has probability at most $\exp(-c_0n^2)$.  An exact test suffices: divide
successively by the integers from $2$ through the integer square root of
the candidate, with no divisor found if and only if it is prime.  The
integer square root is found by binary search.  For the
large prime this takes at most
$2^{\eta n/2+O(1)}\poly(n)$ operations per candidate; the smaller prime
tests are cheaper.  The number of prime calls is polynomial in $n$.

Couple these samplers to ideal unlimited rejection samplers by using the
same fresh bits until a cap is reached.  There are at most $2^{O(n)}$
ordinary sampling calls on any capped execution and polynomially many
prime calls.  These upper bounds depend on $n$ and the capped combinatorial
loops, not on $b$; larger input integers increase the deterministic
arithmetic cost per operation.  The probability that the finite
and ideal executions ever differ is therefore
\begin{equation}
 2^{O(n)}2^{-n^3}+\poly(n)\exp(-c_0n^2)=o(1).
 \label{eq:sampling-coupling}
\end{equation}
The bound also holds for calls chosen adaptively, since each new call has
the stated failure bound conditional on its preceding history.  One-way
randomness causes no difficulty: retain phases and block-table labels
when the algorithm needs them again, and discard other random words
after use.  No earlier random bit is subsequently requested from an
external source.

\paragraph{Complex evaluations and summation.}
\label{par:complex-evaluations}
The fixed-accuracy tests of Sections~\ref{sec:filtering}
and~\ref{sec:checksum} use the same deterministic rational approximations
at every occurrence.  Their comparison results, including equality at a
threshold, are therefore reproducible.  They are distinct from the more
accurate evaluations of the checksum itself.

For those evaluations, first reduce a phase exponent modulo $P$.  Compute
each resulting unit complex number to absolute error at most
$2^{-2n^2}$, with dyadic real and imaginary parts.  All of this is rational
arithmetic of polynomial bit complexity.  For completeness,
\[
 \frac\pi4=\arctan(1/2)+\arctan(1/3)
\]
computes $\pi$ by alternating power series whose remainders decrease
geometrically.  After the angle is reduced to $[0,2\pi]$, the sine and
cosine Taylor series approximate it to $q$ bits with $O(q+1)$ terms:
their remainders are bounded by $7^{r+1}/(r+1)!$, which is less than
$2^{-q}$ for $r$ a sufficiently large constant multiple of $q+1$.
Computing a few additional bits bounds the effect of the angle error and
the final dyadic rounding.  Exact rational evaluation of these truncated
series uses polynomially many bits and operations; even multiplying all
the $O(n)$ dyadic factors without intermediate rounding remains within
a polynomial bit bound.

To make the error accounting explicit, put $\delta=2^{-2n^2}$.
The checksum term has $n$ factors of modulus at most $2$ and one leading
unit phase.  Approximating each to error $\delta$ and telescoping the
product gives error at most
$(n+1)\delta(2+\delta)^n=2^{-2n^2+O(n)}$.  The exact rational weights in
the estimator, including $L_\pi/(P\mu_\pi(k))$, have numerator
and denominator of $O(n)$ bits on every execution that reaches this
calculation, by \eqref{eq:implicit-match-bitlength}; sampled matches
always have positive multiplicity.  Their magnitudes are at most
$2^{O(n)}$.  The number of evaluated terms and all absolute sum bounds
are also $2^{O(n)}$, uniformly over the capped histories.

Round each weighted contribution to a common dyadic grid with $n^2$
fractional bits before adding it.  This avoids multiplying unrelated
denominators during summation.  Exact addition of the resulting dyadics
requires only $n^2+O(n)$ bits, and the total evaluation and rounding error
is at most
\begin{equation}
 2^{-n^2+O(n)}=o(1).
 \label{eq:finite-evaluation-error}
\end{equation}
Division by the power-of-two sample count $N$ preserves the dyadic
representation and adds only $O(n)$ fractional bits.
Choose $n_*$ so that the error in \eqref{eq:finite-evaluation-error} is
at most the $1/16$ arithmetic allowance
already included in the $1/8$ error bound of
Theorem~\ref{thm:checksum}.
The same construction applies to the direct alias sums in
\eqref{eq:subtract-aliases}.  Only polynomially many capped lists can
contribute retained aliases, so their total number is $2^{O(n)}$ even on
an unsuccessful extraction run.  Deduplication can only decrease this
number.  Increase $n_*$ so that the extra error in each direct alias sum
is at most $1/16$.  Together with the $1/8$ checksum accuracy of
Theorem~\ref{thm:checksum}, this leaves
\begin{equation}
 \max_s|\widehat B_s-B_s|<\frac14
 \label{eq:residual-accuracy}
\end{equation}
whenever extraction and checksum estimation succeed.

\subsection{Correctness and the uniform running-time bound}

\begin{proof}[Proof of Theorem~\ref{thm:main}]
First consider an input above the fixed cutoff that satisfies
\eqref{eq:bit-guard}.  Fix any outcome of the tie weights, and hence all
the integers $c_i,t_s$.  The proof divides according to the expectation
of the deficiency on a uniform half; the algorithm never evaluates this
expectation.

If $\E d(V)>0.01004n$, Proposition~\ref{prop:preliminary} returns an
exact decision except with probability $o(1)$.  If instead
\eqref{eq:low-deficiency} holds, Theorem~\ref{thm:aliases} gives complete
alias extraction with probability $1-o(1)$.  Conditional on each fixed
$P$ and transformed instance, Theorem~\ref{thm:checksum} supplies all
checksum estimates with probability $1-o(1)$, uniformly in those data.
The additional sampling loss is $o(1)$ by
\eqref{eq:sampling-coupling}.  These statements do not condition on the
preliminary outcome \textsc{Continue}.  Equivalently, draw the independent
main-procedure randomness even when the preliminary procedure decides
and would not use it.  On the event that the main procedure is successful,
every preliminary continuation is safe, whereas an earlier exact answer
is already correct.  This establishes the required composition without
changing the probability laws in either component theorem.
Thus, for each fixed choice of tie weights, except with probability
$o(1)$ the algorithm either returns the preliminary exact answer or
uses complete aliases and accurate checksums in its final decision.

For a \textsc{No} instance every transformed exact bin is empty, by
Lemma~\ref{lem:isolation}.  On the event just described, if the
preliminary procedure has not already answered, then
\eqref{eq:residual-accuracy} forces every test to return
\textsc{No}.  For a \textsc{Yes} instance, isolation fails with probability
at most $1/12$.  When it succeeds, some transformed exact bin contains
one subset.  Equations~\eqref{eq:exact-bin-checksum}
and~\eqref{eq:residual-accuracy} give $|\widehat B_s|>3/4$ for that bin,
and the algorithm returns \textsc{Yes}.  All the preceding $o(1)$ bounds
are uniform over the guarded inputs and over the fixed tie weights.
Choose the single constant $n_*$ sufficiently large that their total
contribution is at most $1/12$.  The success probability is then at least
$1-1/12-1/12=5/6$, and in particular at least $2/3$.
This cutoff is effective and independent of $b$ and $c$.  The fixed
filtering functions have explicit derivative bounds and positive lower
bounds for their denominators, so the finite-circle approximation has an
effective Riemann-sum error bound.  The remaining requirements use the
displayed tail estimates, positive exponent gaps, and the explicit prime
density bound of Lemma~\ref{lem:prime-collision}.  Taking a larger fixed
integer meets all of them simultaneously.

We next bound work on every random execution, including executions that
abort.  The preliminary list cap has exponent $0.489985$.  The alias
lists and enumerated pairs have exponent $\tau=0.48996$, with
polynomially many repetitions.  The checksum bound is
\[
 2^{(\tau+20\eps)n+O(\sqrt n)}\poly(n).
\]
The term $O(\sqrt n)$ includes the number of patterns, block-table work,
and direct multiplicity computations.  Bounded rejection sampling and
finite arithmetic multiply these costs by a fixed polynomial in
$n+b+2$.  Trial division has exponent at most
$\eta/2=0.25886$ and polynomially many calls, so it is smaller.  Since
$\tau+20\eps=0.4899602<0.489995$ and
$0.489985<0.489995$, increasing the fixed cutoff once more gives fixed
constants $K_0,K$ such that the total work is at most
\begin{equation}
 K_0\,2^{0.489995n}(n+b+2)^K
 \label{eq:all-history-time}
\end{equation}
on every guarded execution.  Input reads, generation of the fixed
approximation meshes, sorting, storage writes, and all random draws are
included.  In particular, no tables depending on $n$ are supplied to the
program in advance.

This operation bound also proves addressability without treating a
multiword integer as a unit-cost value.  Allocate a new word only when it
is written; the number of allocated words cannot exceed
\eqref{eq:all-history-time}.  Its logarithm is at most
\[
 0.489995n+K\log(n+b+2)+O(1).
\]
For sufficiently large $n$, uniformly for every $b\ge1$, the last two
terms are at most $n+b+2$, because $\log x/x\to0$ uniformly for
$x\ge n+3$.  The resulting bound is smaller than
$w=\lceil4(n+b+\log(n+2))\rceil$.  Every memory address therefore fits
in a word.  A polynomial-bit arithmetic operand is stored as as many
words as necessary, and all operations on those words have already been
charged.

It remains to explain the fallback and the quantifiers.  The program can
find $b$ by scanning the input and counting the bits of its largest
integer using shifts, at polynomial cost in $n+b+2$.  The guard is an
integer comparison: $b^{100000}\le2^n$ is equivalent to
\eqref{eq:bit-guard}.  The cutoff is tested first, and a comparison with
this fixed integer may be implemented in finite control.  On a failed
guard or below $n_*$, divide the input
into two halves, enumerate all their subset sums, sort one list, and
search for a complementary sum for each element of the other.  This
includes the empty subsets, is exact on every input, and uses
$2^{n/2}\poly(n+b+2)$ operations.  Each original subset sum has at most
$b+\lceil\log(n+1)\rceil$ bits, so it fits in a word.  The two lists,
their sorting buffers, and their indices need only
$O(2^{\lceil n/2\rceil})$ one-word records; a direct mergesort implementation
fits the specified address space for every $n\ge2$.  No high-precision
complex arithmetic is used in this fallback.  For $n\ge n_*$, the
polynomial storage used by the guard itself fits by the same
$\log(n+b+2)=o(n+b)$ bound as above.  Thus the program terminates and is
correct on inputs outside the asymptotic guarded range as well.

Finally fix any positive integer $c$.  For all sufficiently large $n$,
every $b\le n^c$ satisfies the guard and $n\ge n_*$.  Substituting this
inequality in \eqref{eq:all-history-time} gives
\[
 K_0\,2^{0.489995n}(n+n^c+2)^K
 =O_c(2^{0.49n}),
\]
because the fixed positive gap $0.49-0.489995=0.000005$ absorbs every
fixed polynomial.  The constants in this final comparison may depend
on $c$, but the program, its cutoffs, its numerical coefficients, and
all its sampling rules do not.  This is exactly the uniform
polynomial-bit quantification of Theorem~\ref{thm:main}.
\end{proof}
